\addcontentsline{toc}{subsubsection}{Solution 34. The Ratio--Product Squeeze}

\begin{dlsummary}[title={\bfseries Solution 34. The Ratio--Product Squeeze},breakable]

\begin{proof}
Define
\[
g(x)=\frac{f(x)}x\qquad(x>0).
\]
The first inequality implies
\[
g(x+y)\ge g(x),
\]
so $g$ is nondecreasing. The second inequality becomes
\[
g(xy)\le g(x)g(y). \tag{1}
\]
Taking $x=y=1$ in (1) gives $g(1)\le g(1)^2$. Since $g(1)>0$,
\[
g(1)\ge1. \tag{2}
\]

Let $a>1$ satisfy the composition condition. By monotonicity and (2),
\[
g(a)\ge g(1)\ge1,
\qquad
f(a)=ag(a)\ge a>1.
\]
Therefore $g(f(a))\ge g(a)\ge1$. On the other hand,
\[
f(f(a))=f(a)g(f(a))\le f(a),
\]
so $g(f(a))\le1$. It follows that
\[
g(f(a))=g(a)=g(1)=1.
\]
By monotonicity,
\[
g(x)=1\qquad(1\le x\le a). \tag{3}
\]

Now let $1/a\le x\le1$. From monotonicity and (3), $g(x)\le1$ and
$g(1/x)=1$. Applying (1) to $x$ and $1/x$ gives
\[
1=g(1)\le g(x)g(1/x)=g(x),
\]
so $g(x)=1$. Hence
\[
g(x)=1\qquad(1/a\le x\le a). \tag{4}
\]

Finally, fix any $x>0$. Choose a positive integer $n$ large enough that
$t=x^{1/n}\in[1/a,a]$. By (4), $g(t)=g(1/t)=1$. Repeatedly applying
(1) yields
\[
g(x)=g(t^n)\le g(t)^n=1,
\qquad
g(1/x)=g((1/t)^n)\le g(1/t)^n=1.
\]
But (1) also gives
\[
1=g(1)\le g(x)g(1/x).
\]
Both factors are positive and at most $1$, so they are both equal to $1$.
In particular, $g(x)=1$. Thus $g\equiv1$, and consequently
\[
f(x)=x\qquad(x>0).
\]
\end{proof}

The identity function satisfies both inequalities with equality and
$f(f(a))=f(a)$, so it is the unique solution.

\end{dlsummary}